\chapter{Role of self-attention in a cognitive architecture}
\label{chapter2b}
%

\startcontents[chapters]
\printmyminitoc

\epigraph{\itshape  Intelligence is not only the ability to reason; it is also the ability to find relevant material in memory and to deploy attention when needed.}{-- Daniel Kahneman}

\section{Introduction}
\label{sec:introc2}

Abstract reasoning refers to our ability to analyze information and discover rules to solve arbitrary tasks, and it is  fundamental to general intelligence in human and non-human animals \citep{gentner1997structure,lovett2017modeling}. It is considered a critical component for the development of artificial intelligence (AI) systems and has rapidly started to gain attention. A growing body of literature suggests that current neural architectures exhibit significant limitations in their ability to solve relatively simple visual cognitive tasks in comparison to humans (see ~\citet{ricci37same} for review). 
Given the vast superiority of animals over state-of-the-art AI systems, it makes sense to turn to brain sciences to find inspiration to leverage brain-like mechanisms to improve the ability of  modern deep neural networks  to solve complex visual reasoning tasks. Indeed, a recent human EEG study has shown that attention and memory processes are needed to solve same-different visual reasoning tasks~\citep{alamia2021differential}. \textcolor{black}{This interplay between attention and memory is previously discussed in \citet{buehner2006cognitive,fougnie2008relationship,cochrane2019fluid} emphasizing that a model must learn to perform attention over the memory for reasoning.}

It is thus not surprising that deep neural networks which lack attention and/or memory system fail to robustly solve visual reasoning problems that involve such same-different judgments~\citep{kim2018not}. Recent computer vision works~\citep{messina2021recurrent,vaishnav2021understanding} have provided further computational evidence for the benefits of attention mechanisms in solving a variety of visual reasoning tasks. Interestingly, in both aforementioned studies, a Transformer module was used to implement a form of attention known as self-attention~\citep{cheng2016long,parikh-etal-2016-decomposable}. In such a static module, attention mechanisms are deployed in parallel across an entire visual scene. By contrast, modern cognitive theories of active vision postulate that the visual system explores the environment dynamically  via sequences of attention shifts to select and route task-relevant information to memory. 
Psychophysics experiments \citep{hayhoe2000vision} on overt visual attention have shown that eye movement patterns are driven according to task-dependent routines. %Visual routines are a fundamental concept for active vision as it provides the means to select the specific portions of the scene to perform visual processing operations. Visual routines are thought to be implemented as sequences of elementary operations needed to complete perceptual tasks involving eye movements and processing foveated scene regions. These routines generate internal queries in response to task-specific goals enlisting all the possible locations where the processing focus could shift the attention (see SI section~\ref{si:visualroutine} for more details). The operational procedure of our proposed guided attention fits well with this theoretical framework. 


Inspired by active vision theory, we describe a dynamic attention mechanism, which we call \textit{guided attention}. Our proposed Guided Attention Module for (visual) Reasoning (GAMR) learns to shift attention dynamically, in a task-dependent manner, based on queries internally generated by an LSTM executive controller. Through extensive experiments on the two visual reasoning challenges, the Synthetic Visual Reasoning Test (SVRT) by~\citet{fleuret2011comparing} and the Abstract Reasoning Task (ART) by~\citet{Webb2021EmergentST}, we demonstrate that our neural architecture is capable of learning complex compositions of relational rules in a data-efficient manner and performs better than other state-of-the-art neural architectures for visual reasoning. Using explainability methods, we further characterize the visual strategies leveraged by the model in order to solve representative reasoning tasks. We demonstrate that our model is compositional -- in that it is able to generalize to novel tasks efficiently and learn novel visual routines by re-composing previously learned elementary operations. It also exhibit zero shot generalization ability -by translating knowledge across the tasks sharing similar abstract rules without the need of re-training.


\section{Related Work}
\label{sec:relwork}

Multiple datasets have been used to assess the visual reasoning ability of neural networks. One of the first challenges included the SVRT. Recently introduced Raven style Progressive Matrix datasets, RAVEN \citep{zhang2019raven}, PGM \citet{barrett2018measuring}, focuses on learning nearly seven unique rules and choose one of the eight choices. However, it was found that the dataset was seriously flawed as it was later found that neural architectures could solve tasks by leveraging shortcuts \citep{hu2020hierarchical,spratley2020closer} which were later removed in I-RAVEN~\citep{hu2021stratified}. Prior work~\citep{kim2018not,vaishnav2021understanding,messina2021recurrent} on SVRT studies has focused on the role of attention in solving some of these more challenging tasks. In SVRT, tasks that involve same-different (SD) judgements appear to be significantly harder for neural networks to learn compared to those involving spatial relation judgement (SR) \citep{stabinger201625,Yihe2019ProgramSP,kim2018not} (see \citet{ricci37same} for review). Motivated by neuroscience principles, \citet{vaishnav2021understanding} studied how the addition of feature-based and spatial attention mechanisms differentially affects the learnability of the tasks. These authors found that SVRT tasks could be further taxonomized according to their differential demands for these two types of attention. In another attempt to leverage a Transformer architecture to incorporate attention mechanisms for visual reasoning,  \citet{messina2021recurrent} proposed a recurrent extension of the classic Vision Transformer block (R-ViT). Spatial attention and feedback connections helped the Transformer to learn visual relations better. The authors compared the accuracy of four same-different (SVRT) tasks (tasks \textit{1,5,20,21}) to demonstrate the efficacy of their model. 

With the introduction of Transformer architecture, attention mechanisms started gaining popularity in computer vision. They can either complement \citep{bello2019attention,vaishnav2021understanding,d2021convit} or completely replace existing CNN architectures \citep{ramachandran2019stand,pmlr-v139-touvron21a,Dosovitskiy2020-iq}. Augmenting the attention networks with the convolution architectures helps them explore the best of the both worlds and train relatively faster. In contrast, stand-alone attention architecture takes time to develop similar inductive biases similar to CNN. As initially introduced by \citet{vaswani2017attention}, Transformer uses a self-attention layer, followed by residual connection and layer normalization and a linear projection layer to compute the association between input tokens. We used a similar system where instead of using a self-attention module, in GAMR, feature-based attention vector (an internally generated query) is obtained via an LSTM to guide the attention module to the location essential for the task and we thus call the model as \textit{guided-attention}. 
Since there could be more than one location where the model will attend, we then implemented a memory bank. \textcolor{black}{ A more closely aligned model with the human visual system is proposed by \citet{mnih2014recurrent} -- Recurrent Attention Model (RAM). It learns a saccadic policy over visual images and is trained using reinforcement learning to learn policies.} The Mnih et al system constitutes an example of overt attention. Conversely, GAMR constitutes an example of a covert attention system and assumes a fixed acuity. 


We took inspiration for the memory bank from ESBN~\citep{Webb2021EmergentST}, where mechanisms for variable binding and indirection were introduced in architecture for visual reasoning with the help of external memory. Variable binding is the ability to bind two representations, and indirection is the mechanism involved in retrieving one representation to refer to the other. 
\textcolor{black}{While ESBN was indeed a source of inspiration, we would like to emphasize that GAMR constitutes a substantial improvement over ESBN. First and foremost, ESBN lacks attention. It requires items/objects to be passed serially one by one and hence it cannot solve SVRT or any other multi-object visual reasoning problems. In a sense, the approach taken in ESBN  is to assume an idealized frontend that uses hard attention to perfectly parse a scene into individual objects and then serially pipe them through the architecture. This is where our work makes a substantial contribution by developing an attention front-end (which is soft and not hard) to sequentially attend to relevant features and route them into memory. We tested this template-matching behavior of the ESBN architecture by training it in the presence of Gaussian noise and spatial jittering. It led to a chance-level performance. }
Here, we build on this work and describe an end-to-end trainable model that learns to individuate task-relevant scenes and store their representations in memory to allow the judging of complex relations between objects. Finally, our relational mechanism is inspired by the work in \citet{santoro2017simple} that introduced a plug-and-play model for computing relations between object-like representations in a network. 

\section{Proposed approach}
\label{sec:model}

\begin{figure*}[ht]
\centering
  \includegraphics[width=1\linewidth]{images/chapter2/GAMR_add2.png}
  \caption{Our proposed \textit{GAMR} architecture is composed of three components: an \textit{encoder} module ($f_e$) builds a representation ($z_{img}$) of an image, a \textit{controller}  guides the attention module to dynamically shift attention, and selectively routes task-relevant object representations ($z_t$) to be stored in a memory bank ($M$). The recurrent controller ($f_s$) generates a query vector ($q_{int_t}$) at each time step to guide the next shift of attention based on the current fixation. After a few shifts of attention,  a \textit{reasoning} module ($r_{\theta}$) learns to identify the relationships between objects stored in memory.} 
  \label{fig:model}
\end{figure*}

Our model can be divided into three components: an encoder, a controller, and a relational module (see Fig.~\ref{fig:model} for an overview). In the \textbf{encoder module}, a low dimensional representation ($z_{img}$) for an input image ($x_{in}$) is created. It includes a feature extraction block ($f_e$) which is composed of five convolutional blocks (Figure~\ref{fig:encoder}). The output of the module is denoted as $z_{img}$ $\in \mathcal{R}^{(128,hw)}$ (with $h$ height and  $w$ width). We applied instance normalization (\textit{iNorm}) ~\citep{ulyanov2016instance} over $z_{img}$ before passing it to the controller for further processing without which the network even struggles to learn even simple relations.

\begin{figure}[htb]
\vskip 0.2in
\begin{center}
\includegraphics[width=1.\linewidth]{images/chapter2/conv_esbn.png} \\
\caption{Encoder module ($f_e$) used in \textit{GAMR}. It consists of four convolutional blocks to process input image of 128$\times$128 resolution}
\label{fig:encoder}
\end{center}
\vskip -0.2in
\end{figure}


\textbf{Controller module} is composed of two blocks: a recurrent neural network ($f_s$) which generates an internal query to guide the attention spotlight over the task relevant features. These features are selected with the help of the second block, i.e., guided attention (\textit{GA}) and are sent to the memory bank (\textit{M}).

After $z_{img}$ is built, a \textit{guided-attention} block is used to extracts the relevant visual information from the input image in a top-down manner at each time step ($t$). This block is responsible for generating context vector ($z_t$) to be stored in the memory bank (\textit{M}) along with all the previous context vectors. They are subsequently accessed again later by a reasoning module. This memory bank is inspired by the differentiable memory used in \citet{Webb2021EmergentST}.  

In the guided attention block, an attention vector ($w_t$ $\in \mathcal{R}^{128}$) is obtained by normalizing the addition of encoded feature ($z_{img}$) with internally  generated query ($q_{int_t}$) fetched by $f_s$. This normalized attention vector is used to re-weight the features at every spatial location of $z_{img}$ to generate the context vector $z_t$ $\in \mathcal{R}^{128}$. 
The recurrent controller ($f_s$) uses a Long Short-Term Memory (LSTM) to provide a query vector ($q_{int_t}$ $\in \mathcal{R}^{128}$) in response to a task-specific goal in order to guide attention for the current time step $t$. $f_s$ also independently generates a gate vector ($g$ $\in \mathcal{R}^{128}$) and output vector ($out$ $\in \mathcal{R}^{512}$) with the help of linear layers. 
The gate ($g$) is later used to shift attention to the next task-relevant feature based on the features previously stored in $M$. On the other hand, the decision layer uses the output vector ($out$) to produce the system classification output.  


The \textbf{relational module} is where the reasoning takes place over the context vector ($z_t$) stored in the memory bank ($M$). This module is composed of a \textcolor{black}{two layered} MLP ($r_{\theta}$) which produces a relational vector ($all_{obj}$) \textcolor{black}{similar to the relational network~\citep{santoro2017simple}}. As we will show in section \ref{subsec:learningcomp}, $r_{\theta}$ learns elementary operations associated with basic relational judgments between context vectors ($z_t$) stored in the memory ($M$). It is concatenated with the output ($out$) of the controller ($f_s$) at the last time step ($t=T$) and passed through the decision layer ($f_{\phi}$) to predict the output ($\hat{y}$) for a particular task. We have summarized the steps in Algorithm~\ref{alg:cap}.

\begin{algorithm}[H]
    \caption{Guided Attention Model for (visual) Reasoning (\textit{GAMR}). $LN$ represents layer normalization \citep{ba2016layer} ($||$) indicates the concatenation of two vectors, forming a new vector. \{, \} indicates the concatenation of a matrix and a vector, forming a matrix with one additional row. $\odot$ represents element-wise multiplication and \{. \} represents the product between a scalar and a vector. (h,w) corresponds to the height and width of the feature map obtained from the encoder ($f_e$).}
    \label{alg:cap}

\begin{algorithmic}
    \State \begin{math}{k}_{r_{t=1}} \leftarrow{} 0 \end{math} \Comment{\textcolor{blue}{\begin{math}\in \mathcal{R}^{128}\end{math}}}
    
    \State \begin{math}{h}_{t=1} \leftarrow{} 0\end{math} 
    \Comment{\textcolor{blue}{\begin{math}\in \mathcal{R}^{512}\end{math}}}
    
    \State \begin{math}{M}_{t=1} \leftarrow{} \{\}\end{math}
    \State \begin{math}z_{img} \leftarrow{}
    f_e(x_{in})\end{math} 
    \Comment{\textcolor{blue}{\begin{math}\in \mathcal{R}^{(hw,128)}\end{math}}}
        \For{\texttt{t in 1...T}} 
            \State \begin{math} out,~ g,~ q_{int_t},~ h_t \leftarrow{} f_s(h_{t-1}, ~k_{r_{t-1}}) \end{math}
            \Comment{\textcolor{blue}{\begin{math}~~~~~~out\in \mathcal{R}^{512}, g\in \mathcal{R}^{128}, q_{int_t}\in \mathcal{R}^{128}\end{math}}}
            
            \State \begin{math} w_t \leftarrow{} LN (z_{img} + q_{int_t}.repeat(hw, axis=1))\end{math}
            \Comment{\textcolor{blue}{\begin{math}\in \mathcal{R}^{(hw,128)}\end{math}}}
            
            \State \begin{math} z_t \leftarrow{} (z_{img}~\odot~w_t.sum(axis=1)).sum(axis=1)\end{math}
            \Comment{\textcolor{blue}{\begin{math}\in \mathcal{R}^{128}\end{math}}}
            
            \If{t is 1}
                \State \begin{math} k_{r_{t}} \leftarrow{} 0 \end{math}
            \Else

                \State \begin{math} k_{r_{t}} \leftarrow{} g ~\odot~ M_{t-1}.sum(axis=1) \end{math}
                
            \EndIf
            \State \begin{math} M_{t} \leftarrow{} \{M_{t-1},~ z_{t}\} \end{math}
            \Comment{\textcolor{blue}{\begin{math}\in \mathcal{R}^{(t,128)}\end{math}}}
          \EndFor
    
    \State \begin{math} {all}_{obj} \leftarrow{} r_{{\theta}}(\sum_{i,j=1}^{T} (M_{v_i},~M_{v_j})) \end{math}
    
    \State \begin{math} \hat{y} \leftarrow{} f_{\phi}({all}_{obj}~||~out) \end{math}

\end{algorithmic}
\end{algorithm}

\section{Method}
\label{sec:exp}

\paragraph{Dataset}

The SVRT dataset is composed of \textit{twenty-three} different binary classification challenges, each representing either a single rule or a composition of multiple rules. A complete list of tasks with sample images from each category is shown in Appendix (Fig.~\ref{fig:exampleSD}, \ref{fig:exampleSR}). We formed four different datasets with 0.5k, 1k, 5k, and 10k training samples to train our model. We used unique sets of 4k and 40k samples for validation and test purposes. Classes are balanced for all the analyses. 

We trained the model from scratch for a maximum of 100 epochs with an early stopping criterion of 99\% accuracy on the validation set as followed in \citet{vaishnav2021understanding} using Adam \citep{kingma2014adam} optimizer and a binary cross-entropy loss. We used a hyperparameter optimization framework \textit{Optuna} \citep{optuna_2019} to get the best learning rates and weight decays for these tasks and reports the test accuracy for the models that gave the best validation scores.


\paragraph{Baselines} For the baselines in this dataset, we compared our architecture performance to a Relational Network ($RN$), a popular architecture for reasoning in VQA. The $RN$ uses the same CNN backbone as \textit{GAMR} with feature maps of dimension $\mathcal{R}^{128,hw}$ where $h=8$ and $w=8$. We consider each spatial location of the encoded feature representation as an object (i.e., $N=8\times 8=64$ object representations). We computed all pairwise combinations between all 64 representations using a shared MLP between all the possible pairs (totalling 4096 pairs). These combinations are then averaged and processed through another MLP to compute a relational feature vector ($all_{obj}$) before the final prediction layer ($f_{\phi}$). In a sense, \textit{GAMR} is a special case of an $RN$ network endowed with the ability to attend to a task-relevant subset ($N$ = 4) of these representations with the help of a controller instead of exhaustively computing all 4,096 possible relations -- thus reducing the computing and memory requirements of the architecture very significantly. 

As an additional baseline model, we used 50 layered \textit{ResNet} \citep{he2016deep} and its Transformer-based self-attention network (\textit{Attn-ResNet}) introduced in \citet{vaishnav2021understanding} \textcolor{black}{and follow the training mechanism as defined in the paper}. These have been previously evaluated on SVRT tasks \citep{Funke2021a,vaishnav2021understanding,messina2021solving,messina2021recurrent}. \textit{Attn-ResNet} serves as a powerful baseline because of more free parameters and a self-attention module to compare the proposed active attention component of \textit{GAMR}. In our proposed method, the controller shifts attention head sequentially to individual task-relevant features against a standard self-attention module -- where all task-relevant features are attended to simultaneously. We also evaluated memory based architecture, ESBN~\citep{Webb2021EmergentST} in which we used a similar encoder to that of \textit{GAMR} and pass the images in sequential order with each shape as a single stimulus and the number of time steps as the number of shapes present in the SVRT task. In order to train these models we used images of dimension $128 \times 128$ for architectures such as \textit{RN, ESBN, GAMR} and $256 \times 256$ for \textit{ResNet, Attn-ResNet} (to be consistent with configuration in \citet{vaishnav2021understanding}).  ResNet-50 ($ResNet$) has 23M parameters, Relation Network ($RN$) has 5.1M parameters, ResNet-50 with attention (\textit{Attn-ResNet}) has 24M parameters and \textit{GAMR} \& \textit{ESBN} both have 6.6M parameters. 
\section{Benchmarking the system}
\label{sec:result}

 
\begin{figure*}[ht]
\vspace{-.3cm}
\begin{center}
\includegraphics[width=.47\textwidth]{images/chapter2/2bar_acc500.png} 
\includegraphics[width=.47\textwidth]{images/chapter2/2bar_acc1k.png} \\
\includegraphics[width=.47\textwidth]{images/chapter2/2bar_acc5k.png} 
\includegraphics[width=.47\textwidth]
{images/chapter2/2bar_acc10k.png} 
\caption{Bar plot analysis for the SVRT tasks grouped in same-different ($SD$) and spatially related ($SR$) tasks. We compared the accuracies of five baseline architectures with \textit{GAMR}. We trained these models with .5k, 1k, 5k and 10k samples.}
\label{fig:box}
\end{center}
\end{figure*}

All twenty-three tasks in the SVRT dataset can be broadly divided into two categories: same-different (SD) and spatial relations (SR), based on the relations involved in the tasks. Same-different tasks (\textit{1, 5, 6, 7, 13, 16, 17, 19, 20, 21, 22}) have been found to be harder for neural networks \citep{Ellis2015UnsupervisedLB,kim2018not,stabinger201625,stabinger2021evaluating,Puebla2021.04.06.438551,messina2021recurrent,vaishnav2021understanding} compared to spatial relations tasks (\textit{2, 3, 4, 8, 9, 10, 11, 12, 14, 15, 18, 23}). 

We analyzed an array of architectures and found that, on average, \textit{GAMR} achieves at least 15\% better test accuracy score on $SD$ tasks for 500 samples. In contrast, for $SR$ tasks, the average accuracy has already reached perfection. We find a similar trend for other architectures when trained with different dataset sizes. Overall, RN (GAMR minus attention) and ESBN struggled to solve SVRT tasks even with 10k training samples, pointing towards the lack of an essential component, such as attention. On the other hand, Attn-ResNet architecture demonstrated the second-best performance, which shows its importance in visual reasoning. Results are summarized in Figure~\ref{fig:box}.

\section{Learning Compositionality}
\label{subsec:learningcomp}

Compositionality is the capacity to understand novel combinations from previously known components. While the human brain learns compositionally, deep learning models work on the learning Single task with a Single Model principle. Below, we provide evidence that \textit{GAMR} is capable of harnessing compositionality. We looked for triplets of tasks $(x,y,z)$ such that $z$ would be a composition of tasks $x$ and $y$. We systematically looked for all such available triplets in the SVRT dataset and found three triplets: (\textit{15, 1, 10}), (\textit{18, 16, 10}) and (\textit{21, 19, 25}). We study the ability of the network to learn to compose a new relation with very few training samples, given that it had previously learned the individual rules. We first trained the model with tasks $x$ and $y$ so that the rules are learned with the help of the reasoning module $r_{\theta}$ \textcolor{black}{which is a two-layered MLP}. We expect that the first layer learns the elementary operation over the context vectors stored in the memory block ($M$), while the second layer learns to combine those operations for the tasks $z$. We freeze the model after training with tasks $x$, $y$ and only fine-tune: (i) a layer to learn to combine elementary operations and (ii) a decision layer ($f_{\phi}$) on tasks $z$ with ten samples per category and 100 epochs in total. Results are shown in Figure~\ref{fig:comp}. \textcolor{black}{As our baseline, we trained the model from scratch on task $z$ from a triplet of tasks (\textit{x,y,z}) to show that the model is exploring indeed compositionality. We also ran an additional control experiment for compositionality choosing the random pair of tasks (\textit{x=5, y=17}) such that the rules are not included in tasks (z) \textit{15, 18}, and \textit{21}. When we evaluated the network in this setup, we found the performance of the network to be at the chance level -- aligning with the claim.}


\begin{figure}[ht]
    \begin{center}
      \includegraphics[width=.7\linewidth]{images/chapter2/comp_test3.png}
      \caption{\textbf{Compositionality test}: We train the model with tasks containing specific rules (e.g., task \textit{1} representing same-different discrimination and  task \textit{10} involving identification if the four shapes form a square or not). We show that with its ability to compose already learned rules, \textit{GAMR} can quickly learn with 10 samples per class to adapt to a novel scenario (e.g., task \textit{15} where the rule is to identify if the four shapes forming a square are identical or not).}  
      \label{fig:comp}
    \end{center}
  
\end{figure}

We selected group corresponding to each tasks (\textit{15, 18, 21}) used for composition. Task \textit{15} has four shapes forming a square and are identical. It can be composed of task \textit{1}, helping to identify the same shapes and task \textit{10}, which helps to learn if the four shapes are forming a square. In task \textit{18}, a rule is needed to be learned related to symmetry along the perpendicular bisector of the image. It can be taken as a composition of task \textit{16} which requires learning mirror reflection of the image along the perpendicular bisector of the image and task \textit{10} in which symmetry could be discovered in between 4 shapes (forming a square). At last, we took task \textit{21}, which involves both scaling and rotation between two shapes in an image. As its compositional elements, we designed a variant where there is only rotation and no scaling and represented it with \textit{25} and combined it with another counterpart of \textit{21} where there is scaling and no rotation, i.e., task \textit{19}.


\section{Zero-shot generalization}
\label{sec:zeroshot}

\begin{table}[htbp]
\centering
    \begin{tabular}{cccccc}
\hline
Training            & Test & \multicolumn{3}{c}{Test Accuracy} \\\cline{3-5} 
Task            & Task & GAMR & Attn-ResNet & ResNet \\ \hline
\midrule
\multirow{3}{*}{1}  & 5    & 72.07          & 53.03 &  \textbf{73.04} \\
                    & 15   & \textbf{92.53} & 92.07 &   78.87  \\
                    & 22   & \textbf{84.91} & 80.10 &   67.15  \\\hline
\multirow{3}{*}{5}  & 1    & \textbf{92.64} & 85.73 &   92.28    \\
                    & 15   & \textbf{84.36} & 62.69 &   49.95    \\
                    & 22   & \textbf{76.47} & 55.69 &   50.19    \\\hline
7                   & 22   & \textbf{83.80} & 79.11 &   50.37   \\\hline
21                  & 15   & \textbf{90.53} & 50.00 &   49.76   \\\hline
23                  & 8    & \textbf{85.84} & 58.90 &   59.25   \\ 
\bottomrule
\end{tabular}

\caption{Test accuracy to show if the model learns the correct rules when we train it with a task and test on a different set of SVRT tasks with \textit{GAMR}, Attention with ResNet50 (\textit{Attn-ResNet}) and ResNet-50 (\textit{ResNet}).}
\label{tab:transfer}
\end{table}


We hypothesize that if a model has learned the abstract rule underlying a given task, it should be able to re-use its knowledge of this task on other novel tasks which share a similar rule. To verify that \textit{GAMR} is indeed able to generalize across tasks that share similar rules, we searched for pairs of tasks in SVRT which were composed of at least one common elementary relation \citep{vaishnav2021understanding} between them. For example, in pair (\textit{1}, \textit{22}), task \textit{1} involves the identification of \underline{\textit{two}} similar shapes in category 1 and task \textit{22} involves the identification of \underline{\textit{three}} similar shapes in category 1. In the selected pair, the category that judges the similar rule should belong to the same class (let us say category 1 in the above example) so that we test for the right learnability. We systematically identified a set $x$ of tasks \textit{1, 5, 7, 21, 23} representing elementary relations such as identifying same-different (\textit{1, 5}), grouping (\textit{7}), learning transformation like scaling and rotation (\textit{21}) and learning insideness (\textit{23}). Then we paired them with other tasks sharing similar relations. 
These pairs are task \textit{1} with each of \textit{5, 15 and 22}, task \textit{5} with each of \textit{1, 15 and 22}. Similarly other pairs of tasks are (\textit{7, 22}), (\textit{21, 15}) and (\textit{23, 8}). 
We separately trained the model on the set $x$ and tested the same model on their respective pairs without finetuning further with any samples from the test set (zero-shot classification). 
We observed that \textit{GAMR} could easily generalize from one task to another without re-training. On the contrary, a chance level performance by \textit{ResNet} shows the network's shortcut learning and rote memorization of task-dependent features. In comparison, \textit{GAMR} exhibits far greater abstraction abilities -- demonstrating an ability to comprehend rules in unseen tasks without any training at all. We further explored the strategies learned by \textit{GAMR} using attribution methods for all the tasks. These attribution methods confirm that \textit{GAMR} does indeed use a similar visual routine between the original task for which it was trained and the new task for which it was never trained. Table~\ref{tab:transfer} summarizes these results. 


\section{Ablation Study}
\label{sec:abl}


\paragraph{Benchmarking guided attention}
\label{sec:ga}

\begin{figure}[ht]
  \begin{center}
  \includegraphics[width=.7\linewidth]{images/chapter2/2ablatn.png}
  \caption{We compared the average accuracy over two sub-clusters of SVRT obtained by \textit{GAMR} with its variant when we replaced the guided-attention module with the self-attention (\textit{GAMR-SA}) and when we completely gave away attention and made it a relational reasoning architecture (\textit{GAMR w/o Atn (RN)}).}  
  \label{fig:gam}
  \end{center}
  
\end{figure}

We evaluated our guided-attention module (\textit{GAMR}) and compared it with alternative systems with comparable base-architecture but endowed with self-attention (\textit{With-SA}) or no attention and/or memory (\textit{GAMR w/o Atn (RN)})  over 23 SVRT tasks \textcolor{black}{for the same number of time steps. In \textit{GAMR with-SA}, we add a self-attention layer in the guided attention module and all three input vectors to the attention module are the same ($z_{img}$). Our intuition is that at each time step, the self-attention mechanism should provide the mechanism to learn to attend to different objects in a scene.} As a side note, \textit{GAMR with-SA} turns out to be similar to ARNe~\citep{hahne2019attention} used for solving Raven's tasks. We found that, on average, our Guided Attention Model's relative performance is 11.1\% better than its SA counterpart and 35.6\% than a comparable system lacking attention (or memory) for $SD$ tasks; similarly, relative improvements for $SR$ tasks are 4.5\% and 10.4\% higher. It shows that \textit{GAMR} is  efficient as it yields a higher performance for the same number (1k) of training samples. Results are shown in Figure~\ref{fig:gam}. 

\begin{figure}[htbp]
    \begin{center}
      % 
      \includegraphics[width=1\linewidth]{images/chapter2/abl_ainorm.png} 
      
      \caption{
      \textbf{Ablation studies}: We pruned separate parts of the model, one at a time: controller output ($out$), attention vector ($w_t$), relational vector ($all_{obj}$), feature channel gain factor ($g$) and instance normalization ($iNorm$) and the bar plot show the variation in performance on SD and SR tasks when trained with 1k samples.
      }  
      \label{fig:abl}
    \end{center}
  
\end{figure}

\begin{figure}[htbp]
    \begin{center}
      % 
      \includegraphics[width=1\linewidth]{images/chapter2/timestep.png} 
      
      \caption{
      \textbf{Time steps visualization}: Figure showing shift of attention with each time step in a task-dependent manner. In the first row, the task is to answer if the two shapes are touching each other from the outside. At each time step, the network explores the area where the shapes are touching each other. In other rows, attribution maps show the shifts over different shapes in an image. The controller module for the task in respective rows shifts attention across different shapes at each time step.
      }  
      \label{fig:ablts}
    \end{center}  
\end{figure}

\textit{GAMR} is a complex model with several component, so we now proceed to study what role different components of the proposed architecture play in it's ability to learn reasoning tasks.  
We studied the effect of these components on SD and SR categories. Our lesioning study revealed that \textit{iNorm} plays a vital role in the model's reasoning and generalization capability even for learning simple rules of $SR$ tasks. Normalizing every sample individually helps the model learn the abstract rules involved in task.  We also found that for $SD$ tasks, excluding vector $out$ from the decision-making process is detrimental. The t-SNE plot 
shows that it encodes the independent abstract representation for the SVRT tasks (Figure~\ref{fig:abstract}). We systematically ran an ablation study to show that each of these components are essential to make it an efficient model. In $no\_all_{obj}$, the model takes a decision based on the final outcome ($out$) of the recurrent module ($f_s$); for $no\_w_{t}$, output of the attention block is used to obtain $z_t$ instead after projecting it on $z_{img}$; for $no\_g$, equal weighing is applied to the feature space of the context vectors stored in the memory block. We have summarized the results in Figure~\ref{fig:abl}. We also plot the attribution maps of the model in Figure~\ref{fig:ablts} at each time step and show the way in which the model attends to task-dependent features while learning the rule.

\begin{figure}[ht]
\centering
  \includegraphics[width=.6\linewidth]{images/chapter2/LSTM.png}
  \caption{\textbf{Abstract variable:} t-SNE plot of the output vector (\textit{out}) obtained from the controller ($f_e$) for all 23 SVRT tasks independently. Each cluster can be clearly identified from other clusters representing different relations learned. Tasks are represented as labels with the same colored box around them placed at the mean location of the cluster.}  
  \label{fig:abstract}
\end{figure}



\section{Additional Experiment}
\label{sec:art}

\begin{figure*}[htbp]
\centering
  \includegraphics[width=1\linewidth]{images/chapter2/ESBN_mareo.png}
  \caption{\textbf{ART for \textit{GAMR}:} (a) Same/different discrimination task. (b) Relational match-to-sample task (answer is 2). (c) Distribution-of-three task (answer is 1). (d) Identity rules task (ABA pattern, answer is 3).} 
  \label{fig:esbnGAMR}
\end{figure*}

\paragraph{Dataset} \citet{Webb2021EmergentST}  proposed four visual reasoning tasks (Figure \ref{fig:esbnGAMR}), that we will henceforth refer to as the \textit{Abstract Reasoning Task} (ART):  (1) a same-different (\textit{SD}) discrimination task, (2) a relation match to sample task (\textit{RMTS}), (3) a distribution of three tasks (\textit{Dist3}) and (4) an identity rule task (\textit{ID}). These four tasks utilize shapes from a set of 100 unique Unicode character images \footnote{\url{https://github.com/taylorwwebb/emergent_symbols}}. They are divided into training and test sets into four generalization regimes using different holdout character sets (m = 0, 50, 85, and 95) from 100 characters. We have described training and test samples and different hyperparameters for all four tasks in section~\ref{supp:hyper}.

\paragraph{Baseline models}
As a baseline, we chose the ESBN \citep{Webb2021EmergentST} along with the two other prevalent reasoning architectures, the Transformer \citep{vaswani2017attention} and Relation Network (RN) \citep{santoro2017simple}. These three share a similar encoder backbone as in \textit{GAMR}. \textcolor{black}{We also ran an additional baseline with \textit{ResNet50} to verify if in a multi-object scenario \textit{GAMR} exploring some biases (like visual entropy) which is otherwise not present when segmented images are passed.} In order to make our baselines stronger, we evaluated these models in their natural order, i.e., by passing a single image at a time. We added a random translation (jittering) for the shapes in the area of $\pm 5$ pixels around the center to prevent these architectures from performing template matching. \textcolor{black}{This jittering increases the complexity of the tasks and hence we have to increase the number of time steps from 4 to 6. This caused the differences in results when compared to the implementation in \citet{Webb2021EmergentST}.}  For \textit{GAMR} and \textit{ResNet50}, we present task-relevant images together as a single stimulus Figure~\ref{fig:esbnGAMR} while jittering each shape. We have also added ART results where each image is centered and put together in a single stimulus in Figure~\ref{fig:esbn_result}. In order to make our architecture choose one option from multiple stimuli (\textit{RMTS}: 2, \textit{Dist3} and \textit{ID}: 4), we concatenate the relational vector ($all_{obj}$) for every stimulus and pass them to a linear layer for final decision.  


\begin{figure*}[htbp]
\begin{center}
\includegraphics[width=.95\linewidth]{images/chapter2/GAMR_result_all_ICLR.png} 

\caption{Test accuracy on ART with different holdout sets when the images are $centered$ and compare the accuracy when shapes are $jittered$ in every image. We find that unlike other baselines experiencing a huge drop in performance when shapes are jittered, GAMR is stable. We plot the average accuracy over ten runs on the dataset. $x$ axis corresponds to the four types of tasks, and $y$ represents the average accuracy score. These tasks are as follows: (a) same-different (SD) discrimination task, (b) Relation match to sample task (RMTS); (c) Distribution of three tasks (Dist3); and (d) Identity rule task (ID).}
\label{fig:esbn_result}
\end{center}
\end{figure*}

\paragraph{Results}
We found a near-chance level (50\%) accuracy for all the baseline models and in all the four generalization regimes for the \textit{SD} and \textit{RMTS} tasks (Figure~\ref{fig:esbbnresultall}) which otherwise performed with perfection when images were centered and passed through the same models.
However, our proposed architecture is robust to handle this jittering, as shown in Figure~\ref{fig:esbn_result} where we compare its performance when images are not jittered. For the other two tasks, \textit{Dist3} and \textit{ID}, baseline models performed better than the chance level (25\%). \textit{ESBN} showed an increasing trend in accuracy for progressively easier generalization conditions approaching 0 holdouts. This points toward the fact that the first three shapes in both tasks allow \textit{ESBN} to consider a translation factor while comparing the next three shapes, letting it choose the correct option appropriately. \textit{RN} and \textit{Transformer} consistently struggled to generalize. \textit{ESBN} (memory-based model) performance on SD tasks in both the visual reasoning datasets show that attention is needed for reasoning. 

\begin{figure}[htbp]
\begin{center}
\includegraphics[width=1\linewidth]{images/chapter2/resnetvsgamr_all.png} 
\caption{\textbf{ART}: Comparing the average performance of \textit{GAMR} with other baselines over 10 runs for different holdout values (m = 0, 50, 85, 95).  These models are evaluated on four types of tasks, i.e., Same-Different (SD), Relation match to sample (RMTS), Distribution of 3 (Dist3) and Identity rules (ID).}
\label{fig:esbbnresultall}
\end{center}
\end{figure}

\section{Hyperparameters}

\label{supp:hyper}
\begin{table}[htbp]
\caption{\textbf{ART}: Number of training and test samples used for four different types of tasks.}
\label{tab:artsamples}
\begin{center}
\begin{tabular}{cccccc}
\\\hline
\textbf{Tasks}      &          & m=0    & m=50   & m=85   & m=95   \\ \hline
\multirow{2}{*}{SD} & Training & 18,810 & 4,900  & 420    & 40     \\
                    & Test     & 990    & 4,900  & 10,000 & 10,000 \\ \hline
RMTS                & Training & 10,000 & 10,000 & 10,000 & 480    \\
Dist3               & Training & 10,000 & 10,000 & 10,000 & 360    \\
ID                  & Training & 10,000 & 10,000 & 10,000 & 8,640  \\
                    & Test     & 10,000 & 10,000 & 10,000 & 10,000 \\ \hline
\end{tabular}
\end{center}
\end{table}

\paragraph{Holdout set} For example, holdout \textit{0} represents a generalization regime in which the test sets contain the same characters as those used during training. At the other extreme, in holdout \textit{95}, the training set contains a minimal number of characters, most of which are actually used for tests. Hence, it is necessary to learn the abstract rule in order to generalize to characters in this regime.  

\begin{table}[H]
\caption{\textbf{ART}: For four different tasks number of epochs and learning rates (LR) used to train different architectures.}
\label{tab:arttraining}
\vskip 0.15in
\begin{center}
\resizebox{\textwidth}{!}{%
\begin{tabular}{ccccccccc}
\hline
\textbf{Tasks} & \multicolumn{2}{c}{m=0}                      & \multicolumn{2}{c}{m=50}                     & \multicolumn{2}{c}{m=85}                     & \multicolumn{2}{c}{m=95} \\ \cline{2-9} 
        & \multicolumn{8}{c}{GAMR}                                                                                                                                     \\ \cline{2-9} 
        & Epoch & \multicolumn{1}{c|}{LR} & Epoch & \multicolumn{1}{c|}{LR} & Epoch & \multicolumn{1}{c|}{LR} & Epoch  & LR \\ \hline 
SD      & 50    & \multicolumn{1}{c|}{0.0001}        & 50    & \multicolumn{1}{c|}{0.0005}        & 100   & \multicolumn{1}{c|}{0.0005}        & 200    & 0.001         \\
RMTS    & 50    & \multicolumn{1}{c|}{0.00005}        & 50    & \multicolumn{1}{c|}{0.0001}        & 50    & \multicolumn{1}{c|}{0.0005}        & 300    & 0.0005        \\
Dist3   & 50    & \multicolumn{1}{c|}{0.00005}       & 50    & \multicolumn{1}{c|}{0.0001}        & 50    & \multicolumn{1}{c|}{0.00005}       & 300    & 0.0005        \\
ID      & 50    & \multicolumn{1}{c|}{0.00005}       & 50    & \multicolumn{1}{c|}{0.00005}       & 50    & \multicolumn{1}{c|}{0.0005}        & 100    & 0.0005        \\ \cline{2-9} 
        & \multicolumn{8}{c}{Other baselines}                                                                                                                           \\ \cline{2-9} 
SD      & 50    & \multicolumn{1}{c|}{0.0005}        & 50    & \multicolumn{1}{c|}{0.0005}        & 100   & \multicolumn{1}{c|}{0.0005}        & 200    & 0.0005        \\
RMTS    & 50    & \multicolumn{1}{c|}{0.0005}        & 50    & \multicolumn{1}{c|}{0.0005}        & 50    & \multicolumn{1}{c|}{0.0005}        & 300    & 0.0005        \\
Dist3   & 50    & \multicolumn{1}{c|}{0.0005}        & 50    & \multicolumn{1}{c|}{0.0005}        & 50    & \multicolumn{1}{c|}{0.0005}        & 300    & 0.0005        \\
ID      & 50    & \multicolumn{1}{c|}{0.0005}        & 50    & \multicolumn{1}{c|}{0.0005}        & 50    & \multicolumn{1}{c|}{0.0005}        & 100    & 0.0005        \\ \hline
\end{tabular}}
\end{center}
\end{table}




\section{Conclusion and limitations}
\label{sec:conc}
In this paper, we described a novel Guided Attention Module for (visual) Reasoning (\textit{GAMR}) to bridge the gap between the reasoning abilities of humans and machines. 
Inspired by the cognitive science literature, our module learns to dynamically allocate attention to task-relevant image locations and store relevant information in memory. Our proposed guided-attention mechanism is shown to outperform the self-attention mechanisms commonly used in vision transformers. Our ablation study demonstrated that an interplay between attention and memory was critical to achieving robust abstract visual reasoning. 
Furthermore, we demonstrated that the resulting systems are capable of solving novel tasks efficiently -- by simply rearranging the elemental processing steps to learn the rules without involving any training. 
We demonstrated GAMR's versatility, robustness, and ability to generalize compositionality through an array of experiments. We achieved state-of-the-art accuracy for the two main visual reasoning challenges in the process.  One limitation of the current approach is that it currently only deals with a fixed number of time steps. Training the model with four time steps was sufficient to solve all SVRT tasks efficiently. However, a more flexible approach is needed to allow the model to  automatically allocate a number of time steps according to the computational demand of the task. GAMR is also limited to covert attention unlike biological system where both covert and overt attention are demonstrated. 
